%==============================================================
%  Figura 5.1 - Strategia di riduzione del rischio
%  Adattamento in italiano della Figura 1 della EN ISO 13849-1
%  Disegno vettoriale: i blocchi sono posizionati in modo relativo,
%  quindi per modificarlo basta agire su testi e distanze.
%==============================================================
\begin{tikzpicture}[
  font=\footnotesize,
  node distance=0.9cm,
  fdbox/.style={draw=black!55, fill=black!6, line width=0.4pt, align=center,
                inner sep=3pt, minimum width=6.2cm, text width=5.7cm},
  fdterm/.style={draw=black!55, fill=white, rounded corners=10pt, line width=0.4pt,
                 align=center, inner sep=3pt},
  fddec/.style={diamond, draw=black!55, fill=black!6, line width=0.4pt, align=center,
                inner sep=1pt, aspect=1.9},
  fdline/.style={draw=black!70, line width=0.5pt},
  fdar/.style={fdline, -{Stealth[length=4.5pt,width=3.5pt]}},
  fdgroup/.style={draw=black!70, line width=0.5pt, dashed, rounded corners=10pt},
  lab/.style={font=\footnotesize, inner sep=1.5pt}
]

%--------------------------------------------------------------
% Colonna di sinistra: valutazione e riduzione del rischio
%--------------------------------------------------------------
\node[fdterm, minimum width=2.6cm, minimum height=0.85cm] (inizio) {INIZIO};

\node[fdbox, below=of inizio] (limiti)
  {Determinazione dei limiti\\ della macchina (vedere punto 5.3{\textsuperscript{a}})};

\node[fdbox, below=of limiti] (pericoli)
  {Identificazione dei pericoli\\ (vedere punto 5.4{\textsuperscript{a}} e Allegato B{\textsuperscript{a}})};

\node[fdbox, below=of pericoli] (stima)
  {Stima del rischio\\ (vedere punto 5.5{\textsuperscript{a}})};

\node[fdbox, below=of stima] (ponder)
  {Ponderazione del rischio\\ (vedere punto 5.6{\textsuperscript{a}})};

\node[fddec, below=of ponder, text width=3.4cm] (dec1)
  {Il rischio è stato\\ adeguatamente ridotto?};

\node[fdbox, below=1.7cm of dec1] (riduzione)
  {Processo di riduzione del rischio\\ per il pericolo:\\ 1 mediante progettazione\\ \phantom{1} intrinsecamente sicura\\ 2 mediante ripari e dispositivi\\ \phantom{2} di protezione\\ 3 mediante informazioni per l'uso\\ (vedere Figura 1{\textsuperscript{a}})};

\node[fddec, below=1.4cm of riduzione, text width=4.0cm] (dec3)
  {La misura di protezione\\ scelta dipende da un\\ sistema di comando?};

%--------------------------------------------------------------
% Punti di innesto sulla colonna di sinistra
%--------------------------------------------------------------
\coordinate (j1) at ($(inizio.south)!0.5!(limiti.north)$);   % rientro da ISO 12100
\coordinate (j2) at ($(limiti.south)!0.5!(pericoli.north)$); % ritorno "sì"
\coordinate (j3) at ($(ponder.south)!0.5!(dec1.north)$);     % ritorno "no"
\coordinate (jr) at ($(dec1.south)!0.5!(riduzione.north)$);  % linea di rientro in basso

%--------------------------------------------------------------
% Colonna di destra
%--------------------------------------------------------------
\node[fdterm, text width=6.0cm, minimum height=1.2cm] (iso) at (12.4,0 |- j1)
  {Valutazione del rischio eseguita\\ in conformità alla ISO 12100};

\node[fddec, text width=3.0cm] (dec2) at (12.4,0 |- j3)
  {Sono generati\\ altri pericoli?};

\node[align=center, text width=5.4cm] (nota) at ($(iso.south)!0.5!(dec2.north)+(-2.6,0)$)
  {Questo processo iterativo\\ di riduzione del rischio deve\\ essere eseguito separatamente\\ per ciascun pericolo in ciascuna\\ condizione di utilizzo (compito).};

\node[fdterm, minimum width=2.2cm, minimum height=0.85cm] (fine) at (9.4,0 |- dec1) {FINE};

\node[fdbox, minimum width=6.4cm, text width=5.9cm] (srpcs)
  at ($(11.8,0 |- riduzione.south)+(0,-1.3)$)
  {Processo iterativo di progettazione\\ delle parti dei sistemi di comando\\ legate alla sicurezza (SRP/CS)\\ (vedere Figura 3{\textsuperscript{b}})};

%--------------------------------------------------------------
% Aree tratteggiate
%--------------------------------------------------------------
\node[fdgroup, fit=(limiti)(pericoli)(stima)(ponder)(dec1), inner sep=7pt] (gruppoA) {};
\coordinate (gb1) at ($(riduzione.east)+(0.45,0.35)$);
\coordinate (gb2) at ($(dec3.east)+(0.45,-0.55)$);
\node[fdgroup, fit=(srpcs)(gb1)(gb2), inner sep=7pt] (gruppoB) {};

%--------------------------------------------------------------
% Collegamenti
%--------------------------------------------------------------
\draw[fdar] (inizio.south) -- (limiti.north);
\draw[fdar] (limiti.south) -- (pericoli.north);
\draw[fdar] (pericoli.south) -- (stima.north);
\draw[fdar] (stima.south) -- (ponder.north);
\draw[fdar] (ponder.south) -- (dec1.north);
\draw[fdar] (dec1.south) -- (riduzione.north);
\draw[fdar] (riduzione.south) -- (dec3.north);
\node[lab, left] at ($(dec1.south)!0.35!(riduzione.north)$) {no};

% rientro dalla valutazione del rischio secondo ISO 12100
\draw[fdar] (iso.west) -- (j1);

% uscita verso FINE
\draw[fdar] (dec1.east) -- (fine.west);
\node[lab, above] at ($(dec1.east)!0.5!(fine.west)$) {sì};

% ciclo "sono generati altri pericoli?"
\draw[fdar] (dec2.west) -- (j3);
\node[lab, above] at ($(dec2.west)!0.45!(j3)$) {no};
\draw[fdar] (dec2.north) |- (j2);
\node[lab, above] at ($(j2)!0.45!(dec2.north |- j2)$) {sì};

% ciclo di progettazione delle SRP/CS
\draw[fdar] (dec3.east) -| (srpcs.south);
\node[lab, above] at ($(dec3.east)!0.5!(srpcs.south |- dec3.east)$) {sì};
\draw[fdline] (srpcs.north) |- (15.7,0 |- jr);
\draw[fdline] (dec3.south) -- ++(0,-0.9) -| (15.7,0 |- jr);
\node[lab, right] at ($(dec3.south)+(0,-0.45)$) {no};
\draw[fdar] (dec2.south |- jr) -- (dec2.south);

%--------------------------------------------------------------
% Note e cornice
%--------------------------------------------------------------
\node[lab, anchor=north west] (notaa)
  at ($(limiti.west |- dec3.south)+(0,-1.5)$) {{\textsuperscript{a}}\quad Si riferisce alla ISO 12100:2010};
\node[lab, anchor=north west] (notab)
  at ($(notaa.south west)+(0,-0.15)$) {{\textsuperscript{b}}\quad Si riferisce alla ISO 13849-1};

\draw[draw=black!55, line width=0.5pt]
  ([shift={(-0.45,0.45)}]current bounding box.north west) rectangle
  ([shift={(0.45,-0.45)}]current bounding box.south east);

\end{tikzpicture}
